We study routing over six knowledge operations: answering from parameters, reading persistent memory, retrieving external evidence, writing reusable memory, applying temporary task-time adaptation, and promoting stable updates durably. The core question is not whether any one mechanism works in isolation, but when a controller should choose among them under one objective. This journal-oriented extension keeps the frozen preprint results and adds broader evidence. First, on a frozen recurring-heavy PopQA slice and a broader seeded two-model PopQA run, the router remains near retrieval quality while reducing repeated retrieval and still depends strongly on memory reads, memory writes, and delayed utility. Second, on the bundled FreshQA-style benchmark, the six-action controller still provides the clean updated-knowledge result from the preprint: higher answer quality and lower stale-answer rate than retrieval-only and memory-only baselines with meaningful temporary adaptation, rare guarded consolidation, and rollback. Third, on a public FreshQA-derived benchmark built from weekly public releases, we now find both a regime boundary and a concrete fix: the benchmark is retrieval-dominant, the full router underperforms strong retrieval-style baselines, but a calibrated router that suppresses delayed utility on answer-changing public cases materially improves public-track quality while preserving the positive recurring-memory regime. Together these results support a narrower but stronger claim: multi-timescale routing is useful and externally validated for recurring-memory regimes and controlled update settings, while public changing-answer benchmarks require regime-sensitive delayed-utility calibration rather than uniform future-value routing.
